\chapter{MERN Security}

A MERN stack has attack surface on both ends: the Express API accepts
arbitrary input from the network, and the React app runs in a browser you
do not control. This chapter walks through the concrete threats that matter
in practice and the specific mitigation for each, most of which are single
middleware additions.

\section{Cross-Site Scripting (XSS)}

If user-supplied content (a task title, a comment) is rendered as raw HTML,
an attacker can inject a \texttt{<script>} tag that runs in every other
user's browser who views that content --- stealing cookies, tokens, or
performing actions as them.

React escapes text content by default, so \texttt{\{task.title\}} is safe.
The risk appears only when you deliberately bypass that escaping.

\begin{lstlisting}[language=JavaScript]
// DANGEROUS -- only use with content you fully trust and have sanitized
<div dangerouslySetInnerHTML={{ __html: task.description }} />

// SAFE -- React escapes this automatically
<div>{task.description}</div>
\end{lstlisting}

\begin{warnbox}[WARNING]
If you must render HTML (e.g.\ a rich-text editor's output), sanitize it
server-side or client-side with a library like \texttt{DOMPurify} before
passing it to \texttt{dangerouslySetInnerHTML}. Never trust that "it's just
our own users" makes stored content safe to render unescaped.
\end{warnbox}

\section{Cross-Site Request Forgery (CSRF)}

CSRF tricks a logged-in user's browser into submitting a request to your
API from another site, using the browser's automatically-attached cookies.
It matters specifically because JWTs stored in HttpOnly cookies (Chapter~19)
are sent automatically by the browser on every request to your domain.

Mitigations: set the auth cookie with \texttt{sameSite: "strict"} or
\texttt{"lax"} (Chapter~19), which stops browsers from sending it on
cross-site requests, and only accept state-changing requests
(\texttt{POST}/\texttt{PUT}/\texttt{DELETE}) with a proper CORS
configuration that whitelists your known frontend origin instead of
\texttt{*}.

\begin{lstlisting}[language=JavaScript]
// server.js
import cors from "cors";

app.use(cors({
  origin: process.env.CLIENT_URL, // e.g. "https://taskmanager.app"
  credentials: true,
}));
\end{lstlisting}

\section{NoSQL Injection}

MongoDB queries built from raw request bodies are vulnerable to operator
injection: instead of a string, an attacker sends an object like
\texttt{\{ "\$gt": "" \}} as the \texttt{password} field, which Mongo
interprets as a query operator rather than a literal value --- potentially
bypassing a login check entirely.

\begin{lstlisting}[language=JavaScript]
// Attacker POSTs:
// { "email": "admin@example.com", "password": { "$gt": "" } }
// Without sanitization, User.findOne({ email, password }) can match
// any document where password is greater than an empty string --- i.e. everything.
\end{lstlisting}

Strip \texttt{\$} and \texttt{.} keys from incoming data before it reaches
a query. Combined with the schema-level validation from Chapter~23, this
closes the operator-injection gap that field-type validation alone does
not.

\begin{lstlisting}[language=JavaScript]
// server.js
import mongoSanitize from "express-mongo-sanitize";

app.use(mongoSanitize()); // strips keys starting with "$" or containing "."
\end{lstlisting}

\begin{notebox}[NOTE]
Mongoose schema types already reject an object where a \texttt{String} is
expected in many cases, but relying on that alone is fragile --- explicit
sanitization middleware closes the gap regardless of schema strictness.
\end{notebox}

\section{Brute Force Login Attempts}

Without a limit, an attacker can script thousands of password guesses
against \texttt{/api/auth/login} per minute. Rate limiting caps how many
attempts a single IP can make in a time window.

\begin{lstlisting}[language=JavaScript]
// middleware/rateLimiter.js
import rateLimit from "express-rate-limit";

export const loginLimiter = rateLimit({
  windowMs: 15 * 60 * 1000, // 15 minutes
  max: 10, // 10 attempts per IP per window
  message: { success: false, message: "Too many login attempts. Try again later." },
  standardHeaders: true,
  legacyHeaders: false,
});

// routes/authRoutes.js
router.post("/login", loginLimiter, loginUser);
\end{lstlisting}

\begin{tipbox}[TIP]
Apply a stricter limiter to \texttt{/login} and \texttt{/register}
specifically, and a looser general limiter to the rest of the API. Locking
down every route at login-attempt strictness would throttle legitimate
heavy users of the app.
\end{tipbox}

\section{Credential Theft via XSS-Accessible Tokens}

Storing a JWT in \texttt{localStorage} makes it readable by any JavaScript
running on the page --- including an injected XSS payload. HttpOnly cookies
(Chapter~19) are invisible to \texttt{document.cookie} and to any
JavaScript, so an XSS bug alone cannot steal the token even if it slips
past your escaping.

\section{Insecure Cookies}

An auth cookie sent without \texttt{secure} and \texttt{httpOnly} flags can
be read by client scripts and transmitted over plain HTTP, where it can be
intercepted on the network.

\begin{lstlisting}[language=JavaScript]
res.cookie("token", jwtToken, {
  httpOnly: true,
  secure: process.env.NODE_ENV === "production", // HTTPS only in prod
  sameSite: "strict",
  maxAge: 7 * 24 * 60 * 60 * 1000,
});
\end{lstlisting}

\section{Common HTTP Header Hardening}

\texttt{helmet} sets a collection of security-related HTTP headers (like
disabling MIME sniffing and setting a baseline Content-Security-Policy)
that block several classes of attack with one line.

\begin{lstlisting}[language=JavaScript]
// server.js
import helmet from "helmet";

app.use(helmet());
\end{lstlisting}

\section{Exposed Secrets}

Committing a \texttt{.env} file (or hardcoding a JWT secret, DB URI, or API
key directly in source) puts credentials in Git history forever, even after
deleting the file in a later commit. See Chapter~38 for the full
environment variable setup; the essential rule is that \texttt{.env} is
always in \texttt{.gitignore} and secrets are only read via
\texttt{process.env}.

\section{Malicious File Uploads}

An unrestricted upload endpoint lets an attacker upload an executable
script disguised as an image, or an oversized file meant to exhaust disk
space. Restrict by MIME type and size at the multer/Cloudinary layer, and
never trust the file extension alone.

\begin{lstlisting}[language=JavaScript]
// middleware/upload.js
import multer from "multer";

const upload = multer({
  limits: { fileSize: 2 * 1024 * 1024 }, // 2MB max
  fileFilter: (req, file, cb) => {
    const allowed = ["image/jpeg", "image/png", "image/webp"];
    if (!allowed.includes(file.mimetype)) {
      return cb(new Error("Only JPEG, PNG, or WEBP images are allowed"));
    }
    cb(null, true);
  },
});

export default upload;
\end{lstlisting}

\section{Input Validation}

Every mitigation above assumes untrusted input is validated at the edge
before it reaches business logic. See Chapter~23 for the
\texttt{express-validator}-based validation layer that rejects malformed
requests (wrong types, missing fields, out-of-range values) before a
controller or query ever runs.

\section{Summary Table}

\begin{center}
\begin{tabularx}{\textwidth}{p{2.6cm}Xp{2.6cm}}
\toprule
\textbf{Threat} & \textbf{Mitigation} & \textbf{Where} \\
\midrule
XSS & Rely on React's default escaping; sanitize any HTML before \texttt{dangerouslySetInnerHTML} & Ch.~36 (rendering) \\
CSRF & \texttt{sameSite} cookie flag + CORS origin whitelist & Ch.~19, Ch.~37 \\
NoSQL injection & \texttt{express-mongo-sanitize} to strip \texttt{\$}/\texttt{.} keys, plus schema validation & Ch.~23, Ch.~37 \\
Brute force login & \texttt{express-rate-limit} on \texttt{/auth} routes & Ch.~37 \\
Credential theft (XSS reading tokens) & HttpOnly cookies instead of \texttt{localStorage} & Ch.~19 \\
Insecure cookies & \texttt{httpOnly}, \texttt{secure}, \texttt{sameSite} flags & Ch.~19 \\
Header-based attacks & \texttt{helmet()} middleware & Ch.~37 \\
Exposed secrets & \texttt{.env} + \texttt{.gitignore}, never hardcode & Ch.~38 \\
Malicious uploads & MIME/type + size limits in multer \texttt{fileFilter} & Ch.~37 \\
Malformed/invalid input & \texttt{express-validator} at the route layer & Ch.~23 \\
\bottomrule
\end{tabularx}
\end{center}

\whatwhyhow
{A secure MERN API layers helmet headers, rate limiting, input sanitization/validation, HttpOnly+secure+sameSite cookies, CORS origin restriction, upload restrictions, and secret management, instead of relying on any single defense.}
{Each threat (XSS, CSRF, NoSQL injection, brute force, credential theft, insecure cookies, exposed secrets, malicious uploads) exploits a different gap, so a single fix like "just validate input" leaves the others open.}
{Add \texttt{helmet()} and \texttt{mongoSanitize()} globally, apply \texttt{express-rate-limit} to auth routes, set cookies with \texttt{httpOnly}/\texttt{secure}/\texttt{sameSite}, whitelist CORS origins, restrict upload MIME types/sizes, and keep all secrets in an untracked \texttt{.env}.}
